\section{System Setup and Configuration}
\label{sec:setup}

\subsection{Software stack}

Two distinct runtimes are involved, and conflating them has been a recurring source of confusion.
The UniTime web application is a Java servlet application backed by a relational database; the
tooling built during this phase---the generators, the evaluator and the learning layer---runs
outside the server as standalone Java and Python processes. Table~\ref{tab:stack} lists each
component together with the place it is used and the file in which the version is recorded.

\begin{table}[H]
\centering
\caption{Software stack, with the component each version applies to.}
\label{tab:stack}
\footnotesize
\begin{tabular}{@{}L{3.2cm}L{2.7cm}L{5.0cm}L{4.0cm}@{}}
\toprule
\textbf{Component} & \textbf{Version} & \textbf{Used for} & \textbf{Recorded in} \\
\midrule
Java (JRE)        & 21              & UniTime web application                     & Semester VI report \\
Apache Tomcat     & 9.0             & servlet container for UniTime               & Semester VI report \\
MySQL             & 8.4             & UniTime persistent storage                  & Semester VI report \\
MySQL JDBC driver & ---             & Tomcat\,$\leftrightarrow$\,MySQL connector  & Semester VI report \\
Java (JDK)        & 17 or newer     & \texttt{unitime-rl} module, headless solver & \texttt{rl/pom.xml}, \texttt{rl/SETUP.md} \\
Maven             & 3.9+            & builds \texttt{unitime-rl}                  & \texttt{rl/SETUP.md} \\
Python            & 3.10+           & data pipeline, evaluator, training code     & \texttt{rl/SETUP.md} \\
CPSolver          & 1.4-SNAPSHOT \newline (or 1.4.94) & constraint solver library   & \texttt{rl/pom.xml}, \texttt{rl/SETUP.md} \\
gRPC / protobuf   & 1.84.0 / 4.36.2 & training environment transport              & \texttt{rl/pom.xml} \\
ONNX Runtime      & 1.30.0          & in-process inference of the frozen policy   & \texttt{rl/pom.xml} \\
\bottomrule
\end{tabular}
\end{table}

\subsection{Resolving the Java version question}

The Semester~VI report records Java~21 as the runtime; the project's own presentation material
records Java~17 for the pipeline and learning components. These are not in conflict, and the
repository settles the question precisely. The Maven module declares
\texttt{<maven.compiler.release>17</maven.compiler.release>}, and both \texttt{rl/README.md} and
\texttt{rl/SETUP.md} state the prerequisite as ``JDK~17+''. Java~17 is therefore the \emph{compile
target} of the project's own Java code, not a restriction on the runtime: bytecode targeting
release~17 runs unchanged on Java~21. The UniTime server continues to run on Java~21 with
Tomcat~9.0 and MySQL~8.4, while the headless solver and the training environment are built against
17 and may be executed on either. No component requires a downgrade.

\subsection{CPSolver dependency}

The module's \texttt{pom.xml} depends by default on a locally installed \texttt{1.4-SNAPSHOT} build
of CPSolver, produced from a sibling clone of the upstream repository, because that is the build
used during development. For a machine without that clone, \texttt{rl/SETUP.md} documents
\texttt{-Dcpsolver.version=1.4.94} as a substitute from Maven Central, noting that the local
snapshot is one commit ahead of the published release---a change to an XML-parser security default,
with no API difference.

\subsection{Solver configuration}

The solver configuration used for the pilot was exported from the running UniTime instance and is
kept in the repository as \texttt{rl/config/unitime-coep-default.properties}, so that the headless
runs of Section~\ref{sec:experiments} use the same weights as the server. The values that matter
for the results reported here are listed in Table~\ref{tab:solverconf}.

\begin{table}[H]
\centering
\caption{Solver parameters relevant to the reported results.}
\label{tab:solverconf}
\footnotesize
\begin{tabular}{@{}L{7.4cm}L{2.6cm}L{5.4cm}@{}}
\toprule
\textbf{Property} & \textbf{Value} & \textbf{Effect} \\
\midrule
\texttt{Basic.Mode}                               & \texttt{Initial} & build a new timetable from scratch \\
\texttt{Basic.DisobeyHard}                        & \texttt{false}   & hard constraints are never violated \\
\texttt{Comparator.StudentConflictWeight}         & \texttt{0.2}     & price of each student clash (swept in Section~\ref{sec:experiments}) \\
\texttt{Comparator.HardStudentConflictWeight}     & \texttt{0.8}     & extra price when the student has no alternative section \\
\texttt{Comparator.RoomPreferenceWeight}          & \texttt{1.0}     & weight of room preference violations \\
\texttt{Comparator.TimePreferenceWeight}          & \texttt{0.3}     & weight of time preference violations \\
\texttt{Comparator.TooBigRoomWeight}              & \texttt{0.1}     & penalty for oversized rooms \\
\texttt{General.CBS}                              & \texttt{true}    & Conflict-based Statistics enabled \\
\texttt{General.AutoSameStudents}                 & \texttt{true}    & automatic \texttt{SAME\_STUDENTS} constraints \\
\texttt{General.CompleteSolutionFixInterval}      & \texttt{-1}      & disables the periodic fix pass \\
\texttt{Construction.Class}                       & \multicolumn{2}{L{8.0cm}}{\texttt{NeighbourSelectionWithSuggestions}} \\
\bottomrule
\end{tabular}
\end{table}

Three further settings were changed during this phase and are the precondition for every result in
Section~\ref{sec:experiments}. Their purpose is to stop UniTime from re-sectioning students, so
that the measured conflict count reflects the timetable rather than a reshuffling of cohorts:

\begin{enumerate}
  \item \emph{Student Course Demands} set to \textbf{Enrolled Student Course Demands}, so that the
    solver reads the imported \texttt{StudentClassEnrollment} rows instead of projecting demand from
    curricula. This is a configuration made in
    \textbf{Administration\,$\rightarrow$\,Solver\,$\rightarrow$\,Configurations} and is not stored
    in the exported properties file, which still records the previous value.
  \item \texttt{General.SwitchStudents=false}, which disables the final student-switching pass.
  \item \texttt{Global.LoadStudentEnrlsFromSolution=true}, which loads each student's sections from
    the stored enrolments rather than recomputing them.
\end{enumerate}

With the first setting left at its default of \emph{Projected Student Course Demands}, and with no
last-like enrolment data in the session, the solver loads no students at all, reports
\texttt{Student conflicts: 0}, and exports an empty student block. This failure is silent, and
diagnosing it accounted for a substantial part of the debugging described in
Section~\ref{sec:simulation}.
